\label{sec:lessons}

The first lesson learned is that a literature-backed recommendation is a starting point, not a verdict: it has to be re-tested against the project's own data and constraints, because the conditions that made it the right answer in the source material do not automatically hold here. The disparity-backend comparison (Section~\ref{sec:architecture}) showed that a generic photometric-quality score, the kind of metric usually reported in the stereo-matching literature, was nearly uninformative between every backend that produced disparity at all, while an application-level signal (whether the downstream wave estimator could actually use the resulting point cloud, and how often it had to discard frames) was what actually separated the backends that solved the project's problem from those that did not. The literature pointed in a reasonable direction, but only testing in this project's specific context revealed whether it actually applied.

The second lesson is that fixing a result without understanding why it was wrong tends to produce another wrong result, just a less obviously wrong one. Several diagnoses in Section~\ref{sec:validation} only worked because the investigation went one level past ``the number looks bad'': the altimeter's implausible wave height was traced specifically to a broadband sensor-noise floor outside the wave band, not just patched by clamping the output, which is why band-limiting the spectral integration fixed it without also needing a heave correction that the noise floor had been masking. The same applies to the disparity comparison: StereoBM's apparent total failure under default parameters turned out to be a parameter-choice artifact rather than evidence that block matching cannot work on water at all, and SGM-CUDA's failure to ever support a usable point cloud was traced to the specific GPU binding lacking a left--right consistency control, not treated as a verdict on SGM in general.

The third lesson is that without a validation, comparison, or ground-truth signal, a result is just a number, with no way to tell whether it is measuring the real world or an artifact of the pipeline. No wave buoy or other ground truth was available for this dataset, so every claim in this report rests on cross-checking two estimators that fail in different, independent ways: a temporal, single-point altimeter signal against a spatial, multi-point stereo-geometry signal (Section~\ref{sec:validation}), and the navigation/altitude pipeline against the navigation filter's own modelled water level (Section~\ref{sec:results}). Every wave-estimation bug described in this report, the altimeter noise floor, the point-cloud bad-frame contamination, the height and period mismatches between the two estimators, was only caught because there were two independent numbers to disagree with each other in the first place; any one of them taken alone would have looked self-consistent and been wrong. That mutual agreement after the corrections, not either estimator individually, is what gives the project's wave-parameter outputs a more reliable meaning rather than just plausible-looking output.